\documentclass[a4paper,landscape,8pt]{extarticle}
\usepackage[landscape,margin=0.7cm,top=0.5cm,bottom=0.7cm]{geometry}
\usepackage[utf8]{inputenc}
\usepackage[T1]{fontenc}
\usepackage[ngerman]{babel}
\usepackage{helvet}
\renewcommand{\familydefault}{\sfdefault}
\usepackage{xcolor}
\usepackage{tikz}
\usetikzlibrary{positioning}
\usepackage{enumitem}
\usepackage{multicol}
\usepackage{array}
\usepackage{fancyhdr}
\usepackage{amssymb}
\usepackage{eurosym}

% Farben
\definecolor{bankblue}{HTML}{1565C0}
\definecolor{securitygreen}{HTML}{2E7D32}
\definecolor{warnred}{HTML}{C62828}
\definecolor{lightgray}{HTML}{F5F5F5}
\definecolor{bordergray}{HTML}{BBBBBB}

\pagestyle{fancy}
\fancyhf{}
\renewcommand{\headrulewidth}{0pt}
\fancyfoot[C]{\footnotesize Seite \thepage{} von 3}

\setlength{\parindent}{0pt}
\setlength{\parskip}{2pt}

% Kompakte Listen
\setlist{nosep, leftmargin=1em, itemsep=0pt, topsep=1pt}

\begin{document}

% ============================================================================
% SEITE 1: INFORMATIONSBLATT
% ============================================================================

\noindent{\Large\bfseries Online-Banking: Sicher und einfach} \hfill
{\footnotesize\textbf{Lernziele:} PIN \& TAN verstehen | Phishing erkennen | Sicherheitsregeln}

\vspace{1mm}
\hrule
\vspace{1mm}

% Drei-Spalten-Layout
\begin{minipage}[t]{0.32\linewidth}
\begin{center}
\colorbox{bankblue!15}{%
\begin{minipage}{0.95\linewidth}
\vspace{2mm}
\begin{center}
{\Large\bfseries\textcolor{bankblue}{1. Grundlagen}}\\[1pt]
{\footnotesize\itshape PIN, TAN \& Zwei-Faktor}
\end{center}
\vspace{2mm}
\end{minipage}%
}
\end{center}

\vspace{1mm}

\textbf{Was ist Online-Banking?}\\
Bankgeschäfte vom Handy oder Computer -- wie eine Bankfiliale in der Tasche!

\textbf{Die wichtigsten Begriffe:}
\begin{itemize}
\item \textbf{PIN:} Dein persönlicher ``Schlüssel'' (bleibt gleich)
\item \textbf{TAN:} Einmal-Code für jede Überweisung
\item \textbf{Zwei-Faktor:} Zwei verschiedene Wege zur Bestätigung
\end{itemize}

\textbf{Analogie:}
\begin{itemize}
\item PIN = Schlüssel zum Schließfach
\item TAN = Einmalcode wie beim Paket abholen
\item Zwei-Faktor = Haustür + Tresor (zwei Schlösser)
\end{itemize}

\vspace{2mm}
\fcolorbox{bankblue}{bankblue!5}{%
\begin{minipage}{0.92\linewidth}
\footnotesize
\textbf{Banking-App empfohlen!}\\
Die App deiner Bank ist sicherer und einfacher als die Website.
\end{minipage}%
}

\end{minipage}%
\hfill%
\begin{minipage}[t]{0.32\linewidth}
\begin{center}
\colorbox{warnred!15}{%
\begin{minipage}{0.95\linewidth}
\vspace{2mm}
\begin{center}
{\Large\bfseries\textcolor{warnred}{2. Phishing erkennen}}\\[1pt]
{\footnotesize\itshape Betrüger entlarven}
\end{center}
\vspace{2mm}
\end{minipage}%
}
\end{center}

\vspace{1mm}

\textbf{Was ist Phishing?}\\
Betrüger geben sich als Bank aus, um deine Daten zu stehlen.

\textbf{3 Warnsignale:}
\begin{enumerate}
\item \textbf{Dringlichkeit:} ``Sofort handeln, sonst Konto gesperrt!''
\item \textbf{Links:} ``Klicken Sie hier zur Verifizierung''
\item \textbf{Datenfrage:} ``Nennen Sie Ihre PIN/TAN''
\end{enumerate}

\textbf{Beispiel Phishing-Mail:}\\
{\footnotesize\itshape ``Sehr geehrter Kunde, Ihr Konto wird gesperrt. Klicken Sie HIER um zu bestätigen...''}

\vspace{2mm}
\fcolorbox{warnred}{warnred!5}{%
\begin{minipage}{0.92\linewidth}
\footnotesize
\textbf{Merke:} Deine Bank fragt \textbf{NIEMALS} per E-Mail oder Telefon nach PIN oder TAN!
\end{minipage}%
}

\vspace{2mm}

\textbf{Im Zweifel:}
\begin{itemize}
\item Nicht klicken!
\item Bank anrufen (Nummer von der Karte!)
\item E-Mail löschen
\end{itemize}

\end{minipage}%
\hfill%
\begin{minipage}[t]{0.32\linewidth}
\begin{center}
\colorbox{securitygreen!15}{%
\begin{minipage}{0.95\linewidth}
\vspace{2mm}
\begin{center}
{\Large\bfseries\textcolor{securitygreen}{3. 5 Sicherheitsregeln}}\\[1pt]
{\footnotesize\itshape Unbedingt merken!}
\end{center}
\vspace{2mm}
\end{minipage}%
}
\end{center}

\vspace{1mm}

\textbf{Diese 5 Regeln schützen dich:}

\begin{enumerate}
\item \textbf{Niemals} PIN/TAN am Telefon nennen
\item \textbf{Niemals} auf Bank-Links in E-Mails klicken
\item \textbf{Immer} App nutzen oder Adresse selbst eintippen
\item \textbf{Immer} Kontostand regelmäßig prüfen
\item \textbf{Sofort} Bank anrufen bei Verdacht
\end{enumerate}

\vspace{2mm}

\textbf{Notfall-Nummern:}
\begin{itemize}
\item \textbf{Karte sperren:} 116 116 (kostenlos)
\item \textbf{Deine Bank:} Nummer auf der Karte!
\end{itemize}

\vspace{2mm}
\fcolorbox{securitygreen}{securitygreen!5}{%
\begin{minipage}{0.92\linewidth}
\footnotesize
\textbf{Keine Panik!}\\
Online-Banking ist sicher -- wenn du diese Regeln beachtest.
\end{minipage}%
}

\end{minipage}

\vspace{1mm}

% Zusammenfassung
\begin{center}
\fcolorbox{bordergray}{lightgray}{%
\begin{minipage}{0.98\linewidth}
\centering
\vspace{1.5mm}
\textbf{Zusammenfassung:}
\textcolor{bankblue}{\textbf{PIN}} = Dein Schlüssel (geheim halten!) ~$|$~
\textcolor{bankblue}{\textbf{TAN}} = Einmalcode pro Überweisung ~$|$~
\textcolor{warnred}{\textbf{Phishing}} = Betrug erkennen ~$|$~
\textcolor{securitygreen}{\textbf{5 Regeln}} = Dein Schutz
\vspace{1.5mm}
\end{minipage}%
}
\end{center}

\newpage

% ============================================================================
% SEITE 2: ÜBUNGEN
% ============================================================================

\begin{center}
{\LARGE\bfseries Übungen: Online-Banking sicher nutzen}\\[3pt]
{\normalsize Schau dir Seite 1 an und beantworte die Fragen}
\end{center}

\vspace{3mm}
\hrule
\vspace{4mm}

\begin{multicols}{2}

\textbf{\large Teil A: Grundlagen}

Richtig oder Falsch? Kreuze an.

\vspace{2mm}

\renewcommand{\arraystretch}{1.6}
\begin{tabular}{@{}p{8cm}cc@{}}
& R & F \\
1. Die PIN ist ein Einmalcode für Überweisungen. & $\square$ & $\square$ \\
2. Die TAN gilt nur für eine einzige Überweisung. & $\square$ & $\square$ \\
3. Die Bank fragt manchmal per E-Mail nach der PIN. & $\square$ & $\square$ \\
4. Die Banking-App ist sicherer als der Browser. & $\square$ & $\square$ \\
5. ``Zwei-Faktor'' bedeutet zwei verschiedene Bestätigungen. & $\square$ & $\square$ \\
\end{tabular}
\renewcommand{\arraystretch}{1}

\vspace{6mm}

\textbf{\large Teil B: Phishing erkennen}

Welche 3 Warnsignale deuten auf Phishing hin?

\vspace{2mm}

\begin{enumerate}
\item \dotfill
\item \dotfill
\item \dotfill
\end{enumerate}

\vspace{4mm}

\textbf{Ist diese E-Mail echt oder Phishing?}

\vspace{2mm}

{\footnotesize\itshape ``Sehr geehrter Kunde, wir haben ungewöhnliche Aktivitäten festgestellt. Klicken Sie HIER um Ihr Konto zu bestätigen. Sie haben 24 Stunden Zeit!''}

\vspace{2mm}

$\square$ Echt ~~~~ $\square$ Phishing

\vspace{2mm}

Warum? \dotfill

\columnbreak

\textbf{\large Teil C: Die 5 Sicherheitsregeln}

Vervollständige die Regeln:

\vspace{3mm}

1. \textbf{Niemals} PIN/TAN am \dotfill nennen

\vspace{2mm}

2. \textbf{Niemals} auf \dotfill in E-Mails klicken

\vspace{2mm}

3. \textbf{Immer} App nutzen oder Adresse \dotfill

\vspace{2mm}

4. \textbf{Immer} \dotfill regelmäßig prüfen

\vspace{2mm}

5. \textbf{Sofort} Bank \dotfill bei Verdacht

\vspace{6mm}

\textbf{\large Teil D: Notfall-Nummern}

Fülle aus:

\vspace{3mm}

Karte sperren (bundesweit): \dotfill

\vspace{3mm}

Wo steht die Nummer deiner Bank? \dotfill

\vspace{6mm}

\textbf{\large Teil E: Analogien verstehen}

Verbinde:

\vspace{2mm}

\renewcommand{\arraystretch}{1.8}
\begin{tabular}{@{}p{3cm}p{4.5cm}@{}}
PIN & $\circ$ ~~~~ $\circ$ Einmalcode beim Paket \\
TAN & $\circ$ ~~~~ $\circ$ Schlüssel zum Schließfach \\
Zwei-Faktor & $\circ$ ~~~~ $\circ$ Haustür + Tresor \\
\end{tabular}
\renewcommand{\arraystretch}{1}

\end{multicols}

\newpage

% ============================================================================
% SEITE 3: CHECKLISTE
% ============================================================================

\begin{center}
{\LARGE\bfseries Checkliste: Das weiß ich jetzt}\\[3pt]
{\normalsize Geh die Liste durch und hake ab, was du sicher verstanden hast}
\end{center}

\vspace{3mm}
\hrule
\vspace{6mm}

% Drei Boxen nebeneinander
\begin{minipage}[t]{0.31\linewidth}
\begin{center}
\fcolorbox{bankblue}{bankblue!10}{%
\begin{minipage}{0.92\linewidth}
\vspace{3mm}
\begin{center}
{\large\bfseries\textcolor{bankblue}{Grundlagen}}
\end{center}
\vspace{2mm}
\begin{itemize}[label=$\square$, itemsep=5pt]
\item Ich weiß, was \textbf{PIN} ist
\item Ich weiß, was \textbf{TAN} ist
\item Ich verstehe \textbf{Zwei-Faktor}
\item Banking-App ist \textbf{sicherer}
\end{itemize}
\vspace{3mm}
\end{minipage}%
}
\end{center}
\end{minipage}%
\hfill%
\begin{minipage}[t]{0.31\linewidth}
\begin{center}
\fcolorbox{warnred}{warnred!10}{%
\begin{minipage}{0.92\linewidth}
\vspace{3mm}
\begin{center}
{\large\bfseries\textcolor{warnred}{Phishing erkennen}}
\end{center}
\vspace{2mm}
\begin{itemize}[label=$\square$, itemsep=5pt]
\item Ich kenne die \textbf{3 Warnsignale}
\item Ich weiß: Bank fragt \textbf{nie} nach PIN
\item Im Zweifel: \textbf{Nicht klicken}
\item Verdacht: \textbf{Bank anrufen}
\end{itemize}
\vspace{3mm}
\end{minipage}%
}
\end{center}
\end{minipage}%
\hfill%
\begin{minipage}[t]{0.31\linewidth}
\begin{center}
\fcolorbox{securitygreen}{securitygreen!10}{%
\begin{minipage}{0.92\linewidth}
\vspace{3mm}
\begin{center}
{\large\bfseries\textcolor{securitygreen}{Sicherheitsregeln}}
\end{center}
\vspace{2mm}
\begin{itemize}[label=$\square$, itemsep=5pt]
\item Ich kenne die \textbf{5 Regeln}
\item Ich habe die \textbf{Notfallnummer}
\item Ich prüfe regelmäßig den \textbf{Kontostand}
\item Ich fühle mich \textbf{sicherer}
\end{itemize}
\vspace{3mm}
\end{minipage}%
}
\end{center}
\end{minipage}

\vspace{8mm}

% Gesamtverständnis
\begin{center}
\fcolorbox{black}{lightgray}{%
\begin{minipage}{0.95\linewidth}
\vspace{4mm}
\begin{center}
{\large\bfseries Das Wichtigste}
\end{center}
\vspace{3mm}
\begin{multicols}{2}
\begin{itemize}[label=$\square$, itemsep=6pt]
\item Online-Banking ist \textbf{sicher} -- mit den richtigen Regeln
\item Meine Bank fragt \textbf{nie} nach PIN/TAN per E-Mail
\item Bei Verdacht: \textbf{Nicht klicken, Bank anrufen}
\item Die \textbf{Banking-App} ist mein sicherer Zugang
\item Ich kenne die \textbf{Sperr-Nummer} 116 116
\end{itemize}
\end{multicols}
\vspace{3mm}
\end{minipage}%
}
\end{center}

\vspace{8mm}

% Notfall-Notizen
\begin{center}
\fcolorbox{bordergray}{white}{%
\begin{minipage}{0.95\linewidth}
\vspace{4mm}
\begin{center}
{\large\bfseries Meine Banking-Notizen}\\[3pt]
{\footnotesize (Trage hier deine wichtigen Informationen ein -- NICHT die PIN!)}
\end{center}
\vspace{4mm}
\renewcommand{\arraystretch}{2}
\begin{tabular}{@{}p{0.45\linewidth}p{0.45\linewidth}@{}}
\textbf{Meine Bank:} & \textbf{Notfall-Nummern:} \\
\dotfill & Karte sperren: 116 116 \\[2mm]
\textbf{Telefon meiner Bank:} & \textbf{Banking-App installiert?} \\
\dotfill & $\square$ Ja ~~~ $\square$ Noch nicht \\[2mm]
\end{tabular}
\renewcommand{\arraystretch}{1}
\vspace{4mm}
\end{minipage}%
}
\end{center}

\vfill

\begin{center}
\footnotesize\textit{Stand: \today{} -- Online-Banking ist sicher, wenn du die 5 Regeln beachtest. Deine Bank fragt NIE per E-Mail nach PIN/TAN!}
\end{center}

\end{document}
